\yoschapterstar{Contents}\label{chap:toc}
\vspace{2.6mm}\noindent\hangindent=1.6em\hangafter=1\textbf{0.\ \ Preliminaries}\nobreak\hfill\nobreak\leaders\hbox to 0.55em{\hss.\hss}\hfill\nobreak\ \pageref{chap:0}\par
\noindent\hangindent=3.5em\hangafter=1\hspace*{3.5em}1.\ Set Theory\nobreak\hfill\nobreak\leaders\hbox to 0.55em{\hss.\hss}\hfill\nobreak\ \pageref{sec:0.1}\par
\noindent\hangindent=3.5em\hangafter=1\hspace*{3.5em}2.\ Topological Spaces\nobreak\hfill\nobreak\leaders\hbox to 0.55em{\hss.\hss}\hfill\nobreak\ \pageref{sec:0.2}\par
\noindent\hangindent=3.5em\hangafter=1\hspace*{3.5em}3.\ Measure Spaces\nobreak\hfill\nobreak\leaders\hbox to 0.55em{\hss.\hss}\hfill\nobreak\ \pageref{sec:0.3}\par
\noindent\hangindent=3.5em\hangafter=1\hspace*{3.5em}4.\ Linear Spaces\nobreak\hfill\nobreak\leaders\hbox to 0.55em{\hss.\hss}\hfill\nobreak\ \pageref{sec:0.4}\par
\vspace{2.6mm}\noindent\hangindent=1.6em\hangafter=1\textbf{I.\ \ Semi-norms}\nobreak\hfill\nobreak\leaders\hbox to 0.55em{\hss.\hss}\hfill\nobreak\ \pageref{chap:I}\par
\noindent\hangindent=3.5em\hangafter=1\hspace*{3.5em}1.\ Semi-norms and Locally Convex Linear Topological Spaces\nobreak\hfill\nobreak\leaders\hbox to 0.55em{\hss.\hss}\hfill\nobreak\ \pageref{sec:I.1}\par
\noindent\hangindent=3.5em\hangafter=1\hspace*{3.5em}2.\ Norms and Quasi-norms\nobreak\hfill\nobreak\leaders\hbox to 0.55em{\hss.\hss}\hfill\nobreak\ \pageref{sec:I.2}\par
\noindent\hangindent=3.5em\hangafter=1\hspace*{3.5em}3.\ Examples of Normed Linear Spaces\nobreak\hfill\nobreak\leaders\hbox to 0.55em{\hss.\hss}\hfill\nobreak\ \pageref{sec:I.3}\par
\noindent\hangindent=3.5em\hangafter=1\hspace*{3.5em}4.\ Examples of Quasi-normed Linear Spaces\nobreak\hfill\nobreak\leaders\hbox to 0.55em{\hss.\hss}\hfill\nobreak\ \pageref{sec:I.4}\par
\noindent\hangindent=3.5em\hangafter=1\hspace*{3.5em}5.\ Pre-Hilbert Spaces\nobreak\hfill\nobreak\leaders\hbox to 0.55em{\hss.\hss}\hfill\nobreak\ \pageref{sec:I.5}\par
\noindent\hangindent=3.5em\hangafter=1\hspace*{3.5em}6.\ Continuity of Linear Operators\nobreak\hfill\nobreak\leaders\hbox to 0.55em{\hss.\hss}\hfill\nobreak\ \pageref{sec:I.6}\par
\noindent\hangindent=3.5em\hangafter=1\hspace*{3.5em}7.\ Bounded Sets and Bornologic Spaces\nobreak\hfill\nobreak\leaders\hbox to 0.55em{\hss.\hss}\hfill\nobreak\ \pageref{sec:I.7}\par
\noindent\hangindent=3.5em\hangafter=1\hspace*{3.5em}8.\ Generalized Functions and Generalized Derivatives\nobreak\hfill\nobreak\leaders\hbox to 0.55em{\hss.\hss}\hfill\nobreak\ \pageref{sec:I.8}\par
\noindent\hangindent=3.5em\hangafter=1\hspace*{3.5em}9.\ $B$-spaces and $F$-spaces\nobreak\hfill\nobreak\leaders\hbox to 0.55em{\hss.\hss}\hfill\nobreak\ \pageref{sec:I.9}\par
\noindent\hangindent=3.5em\hangafter=1\hspace*{3.5em}10.\ The Completion\nobreak\hfill\nobreak\leaders\hbox to 0.55em{\hss.\hss}\hfill\nobreak\ \pageref{sec:I.10}\par
\noindent\hangindent=3.5em\hangafter=1\hspace*{3.5em}11.\ Factor Spaces of a $B$-space\nobreak\hfill\nobreak\leaders\hbox to 0.55em{\hss.\hss}\hfill\nobreak\ \pageref{sec:I.11}\par
\noindent\hangindent=3.5em\hangafter=1\hspace*{3.5em}12.\ The Partition of Unity\nobreak\hfill\nobreak\leaders\hbox to 0.55em{\hss.\hss}\hfill\nobreak\ \pageref{sec:I.12}\par
\noindent\hangindent=3.5em\hangafter=1\hspace*{3.5em}13.\ Generalized Functions with Compact Support\nobreak\hfill\nobreak\leaders\hbox to 0.55em{\hss.\hss}\hfill\nobreak\ \pageref{sec:I.13}\par
\noindent\hangindent=3.5em\hangafter=1\hspace*{3.5em}14.\ The Direct Product of Generalized Functions\nobreak\hfill\nobreak\leaders\hbox to 0.55em{\hss.\hss}\hfill\nobreak\ \pageref{sec:I.14}\par
\vspace{2.6mm}\noindent\hangindent=1.6em\hangafter=1\textbf{II.\ \ Applications of the Baire-Hausdorff Theorem}\nobreak\hfill\nobreak\leaders\hbox to 0.55em{\hss.\hss}\hfill\nobreak\ \pageref{chap:II}\par
\noindent\hangindent=3.5em\hangafter=1\hspace*{3.5em}1.\ The Uniform Boundedness Theorem and the Resonance Theorem\nobreak\hfill\nobreak\leaders\hbox to 0.55em{\hss.\hss}\hfill\nobreak\ \pageref{sec:II.1}\par
\noindent\hangindent=3.5em\hangafter=1\hspace*{3.5em}2.\ The Vitali-Hahn-Saks Theorem\nobreak\hfill\nobreak\leaders\hbox to 0.55em{\hss.\hss}\hfill\nobreak\ \pageref{sec:II.2}\par
\noindent\hangindent=3.5em\hangafter=1\hspace*{3.5em}3.\ The Termwise Differentiability of a Sequence of Generalized Functions\nobreak\hfill\nobreak\leaders\hbox to 0.55em{\hss.\hss}\hfill\nobreak\ \pageref{sec:II.3}\par
\noindent\hangindent=3.5em\hangafter=1\hspace*{3.5em}4.\ The Principle of the Condensation of Singularities\nobreak\hfill\nobreak\leaders\hbox to 0.55em{\hss.\hss}\hfill\nobreak\ \pageref{sec:II.4}\par
\noindent\hangindent=3.5em\hangafter=1\hspace*{3.5em}5.\ The Open Mapping Theorem\nobreak\hfill\nobreak\leaders\hbox to 0.55em{\hss.\hss}\hfill\nobreak\ \pageref{sec:II.5}\par
\noindent\hangindent=3.5em\hangafter=1\hspace*{3.5em}6.\ The Closed Graph Theorem\nobreak\hfill\nobreak\leaders\hbox to 0.55em{\hss.\hss}\hfill\nobreak\ \pageref{sec:II.6}\par
\noindent\hangindent=3.5em\hangafter=1\hspace*{3.5em}7.\ An Application of the Closed Graph Theorem (Hörmander's Theorem)\nobreak\hfill\nobreak\leaders\hbox to 0.55em{\hss.\hss}\hfill\nobreak\ \pageref{sec:II.7}\par
\vspace{2.6mm}\noindent\hangindent=1.6em\hangafter=1\textbf{III.\ \ The Orthogonal Projection and F. Riesz' Representation Theorem}\nobreak\hfill\nobreak\leaders\hbox to 0.55em{\hss.\hss}\hfill\nobreak\ \pageref{chap:III}\par
\noindent\hangindent=3.5em\hangafter=1\hspace*{3.5em}1.\ The Orthogonal Projection\nobreak\hfill\nobreak\leaders\hbox to 0.55em{\hss.\hss}\hfill\nobreak\ \pageref{sec:III.1}\par
\noindent\hangindent=3.5em\hangafter=1\hspace*{3.5em}2.\ "Nearly Orthogonal" Elements\nobreak\hfill\nobreak\leaders\hbox to 0.55em{\hss.\hss}\hfill\nobreak\ \pageref{sec:III.2}\par
\noindent\hangindent=3.5em\hangafter=1\hspace*{3.5em}3.\ The Ascoli-Arzelà Theorem\nobreak\hfill\nobreak\leaders\hbox to 0.55em{\hss.\hss}\hfill\nobreak\ \pageref{sec:III.3}\par
\noindent\hangindent=3.5em\hangafter=1\hspace*{3.5em}4.\ The Orthogonal Base. Bessel's Inequality and Parseval's Relation\nobreak\hfill\nobreak\leaders\hbox to 0.55em{\hss.\hss}\hfill\nobreak\ \pageref{sec:III.4}\par
\noindent\hangindent=3.5em\hangafter=1\hspace*{3.5em}5.\ E. \textsc{Schmidt}'s Orthogonalization\nobreak\hfill\nobreak\leaders\hbox to 0.55em{\hss.\hss}\hfill\nobreak\ \pageref{sec:III.5}\par
\noindent\hangindent=3.5em\hangafter=1\hspace*{3.5em}6.\ F. \textsc{Riesz}' Representation Theorem\nobreak\hfill\nobreak\leaders\hbox to 0.55em{\hss.\hss}\hfill\nobreak\ \pageref{sec:III.6}\par
\noindent\hangindent=3.5em\hangafter=1\hspace*{3.5em}7.\ The Lax-Milgram Theorem\nobreak\hfill\nobreak\leaders\hbox to 0.55em{\hss.\hss}\hfill\nobreak\ \pageref{sec:III.7}\par
\noindent\hangindent=3.5em\hangafter=1\hspace*{3.5em}8.\ A Proof of the Lebesgue-Nikodym Theorem\nobreak\hfill\nobreak\leaders\hbox to 0.55em{\hss.\hss}\hfill\nobreak\ \pageref{sec:III.8}\par
\noindent\hangindent=3.5em\hangafter=1\hspace*{3.5em}9.\ The Aronszajn-Bergman Reproducing Kernel\nobreak\hfill\nobreak\leaders\hbox to 0.55em{\hss.\hss}\hfill\nobreak\ \pageref{sec:III.9}\par
\noindent\hangindent=3.5em\hangafter=1\hspace*{3.5em}10.\ The Negative Norm of P. \textsc{Lax}\nobreak\hfill\nobreak\leaders\hbox to 0.55em{\hss.\hss}\hfill\nobreak\ \pageref{sec:III.10}\par
\noindent\hangindent=3.5em\hangafter=1\hspace*{3.5em}11.\ Local Structures of Generalized Functions\nobreak\hfill\nobreak\leaders\hbox to 0.55em{\hss.\hss}\hfill\nobreak\ \pageref{sec:III.11}\par
\vspace{2.6mm}\noindent\hangindent=1.6em\hangafter=1\textbf{IV.\ \ The Hahn-Banach Theorems}\nobreak\hfill\nobreak\leaders\hbox to 0.55em{\hss.\hss}\hfill\nobreak\ \pageref{chap:IV}\par
\noindent\hangindent=3.5em\hangafter=1\hspace*{3.5em}1.\ The Hahn-Banach Extension Theorem in Real Linear Spaces\nobreak\hfill\nobreak\leaders\hbox to 0.55em{\hss.\hss}\hfill\nobreak\ \pageref{sec:IV.1}\par
\noindent\hangindent=3.5em\hangafter=1\hspace*{3.5em}2.\ The Generalized Limit\nobreak\hfill\nobreak\leaders\hbox to 0.55em{\hss.\hss}\hfill\nobreak\ \pageref{sec:IV.2}\par
\noindent\hangindent=3.5em\hangafter=1\hspace*{3.5em}3.\ Locally Convex, Complete Linear Topological Spaces\nobreak\hfill\nobreak\leaders\hbox to 0.55em{\hss.\hss}\hfill\nobreak\ \pageref{sec:IV.3}\par
\noindent\hangindent=3.5em\hangafter=1\hspace*{3.5em}4.\ The Hahn-Banach Extension Theorem in Complex Linear Spaces\nobreak\hfill\nobreak\leaders\hbox to 0.55em{\hss.\hss}\hfill\nobreak\ \pageref{sec:IV.4}\par
\noindent\hangindent=3.5em\hangafter=1\hspace*{3.5em}5.\ The Hahn-Banach Extension Theorem in Normed Linear Spaces\nobreak\hfill\nobreak\leaders\hbox to 0.55em{\hss.\hss}\hfill\nobreak\ \pageref{sec:IV.5}\par
\noindent\hangindent=3.5em\hangafter=1\hspace*{3.5em}6.\ The Existence of Non-trivial Continuous Linear Functionals\nobreak\hfill\nobreak\leaders\hbox to 0.55em{\hss.\hss}\hfill\nobreak\ \pageref{sec:IV.6}\par
\noindent\hangindent=3.5em\hangafter=1\hspace*{3.5em}7.\ Topologies of Linear Maps\nobreak\hfill\nobreak\leaders\hbox to 0.55em{\hss.\hss}\hfill\nobreak\ \pageref{sec:IV.7}\par
\noindent\hangindent=3.5em\hangafter=1\hspace*{3.5em}8.\ The Embedding of $X$ in its Bidual Space $X''$\nobreak\hfill\nobreak\leaders\hbox to 0.55em{\hss.\hss}\hfill\nobreak\ \pageref{sec:IV.8}\par
\noindent\hangindent=3.5em\hangafter=1\hspace*{3.5em}9.\ Examples of Dual Spaces\nobreak\hfill\nobreak\leaders\hbox to 0.55em{\hss.\hss}\hfill\nobreak\ \pageref{sec:IV.9}\par
\vspace{2.6mm}\noindent\hangindent=1.6em\hangafter=1\textbf{V.\ \ Strong Convergence and Weak Convergence}\nobreak\hfill\nobreak\leaders\hbox to 0.55em{\hss.\hss}\hfill\nobreak\ \pageref{chap:V}\par
\noindent\hangindent=3.5em\hangafter=1\hspace*{3.5em}1.\ The Weak Convergence and The Weak* Convergence\nobreak\hfill\nobreak\leaders\hbox to 0.55em{\hss.\hss}\hfill\nobreak\ \pageref{sec:V.1}\par
\noindent\hangindent=3.5em\hangafter=1\hspace*{3.5em}2.\ The Local Sequential Weak Compactness of Reflexive $B$-spaces. The Uniform Convexity\nobreak\hfill\nobreak\leaders\hbox to 0.55em{\hss.\hss}\hfill\nobreak\ \pageref{sec:V.2}\par
\noindent\hangindent=3.5em\hangafter=1\hspace*{3.5em}3.\ Dunford's Theorem and The Gelfand-Mazur Theorem\nobreak\hfill\nobreak\leaders\hbox to 0.55em{\hss.\hss}\hfill\nobreak\ \pageref{sec:V.3}\par
\noindent\hangindent=3.5em\hangafter=1\hspace*{3.5em}4.\ The Weak and Strong Measurability. Pettis' Theorem\nobreak\hfill\nobreak\leaders\hbox to 0.55em{\hss.\hss}\hfill\nobreak\ \pageref{sec:V.4}\par
\noindent\hangindent=3.5em\hangafter=1\hspace*{3.5em}5.\ Bochner's Integral\nobreak\hfill\nobreak\leaders\hbox to 0.55em{\hss.\hss}\hfill\nobreak\ \pageref{sec:V.5}\par
\vspace{2.6mm}\noindent\hangindent=1.6em\hangafter=1\textbf{Appendix to Chapter V. Weak Topologies and Duality in Locally Convex Linear Topological Spaces}\nobreak\hfill\nobreak\leaders\hbox to 0.55em{\hss.\hss}\hfill\nobreak\ \pageref{chap:AV}\par
\noindent\hangindent=3.5em\hangafter=1\hspace*{3.5em}1.\ Polar Sets\nobreak\hfill\nobreak\leaders\hbox to 0.55em{\hss.\hss}\hfill\nobreak\ \pageref{sec:AV.1}\par
\noindent\hangindent=3.5em\hangafter=1\hspace*{3.5em}2.\ Barrel Spaces\nobreak\hfill\nobreak\leaders\hbox to 0.55em{\hss.\hss}\hfill\nobreak\ \pageref{sec:AV.2}\par
\noindent\hangindent=3.5em\hangafter=1\hspace*{3.5em}3.\ Semi-reflexivity and Reflexivity\nobreak\hfill\nobreak\leaders\hbox to 0.55em{\hss.\hss}\hfill\nobreak\ \pageref{sec:AV.3}\par
\noindent\hangindent=3.5em\hangafter=1\hspace*{3.5em}4.\ The Eberlein-Shmulyan Theorem\nobreak\hfill\nobreak\leaders\hbox to 0.55em{\hss.\hss}\hfill\nobreak\ \pageref{sec:AV.4}\par
\vspace{2.6mm}\noindent\hangindent=1.6em\hangafter=1\textbf{VI.\ \ Fourier Transform and Differential Equations}\nobreak\hfill\nobreak\leaders\hbox to 0.55em{\hss.\hss}\hfill\nobreak\ \pageref{chap:VI}\par
\noindent\hangindent=3.5em\hangafter=1\hspace*{3.5em}1.\ The Fourier Transform of Rapidly Decreasing Functions\nobreak\hfill\nobreak\leaders\hbox to 0.55em{\hss.\hss}\hfill\nobreak\ \pageref{sec:VI.1}\par
\noindent\hangindent=3.5em\hangafter=1\hspace*{3.5em}2.\ The Fourier Transform of Tempered Distributions\nobreak\hfill\nobreak\leaders\hbox to 0.55em{\hss.\hss}\hfill\nobreak\ \pageref{sec:VI.2}\par
\noindent\hangindent=3.5em\hangafter=1\hspace*{3.5em}3.\ Convolutions\nobreak\hfill\nobreak\leaders\hbox to 0.55em{\hss.\hss}\hfill\nobreak\ \pageref{sec:VI.3}\par
\noindent\hangindent=3.5em\hangafter=1\hspace*{3.5em}4.\ The Paley-Wiener Theorems. The One-sided Laplace Transform\nobreak\hfill\nobreak\leaders\hbox to 0.55em{\hss.\hss}\hfill\nobreak\ \pageref{sec:VI.4}\par
\noindent\hangindent=3.5em\hangafter=1\hspace*{3.5em}5.\ Titchmarsh's Theorem\nobreak\hfill\nobreak\leaders\hbox to 0.55em{\hss.\hss}\hfill\nobreak\ \pageref{sec:VI.5}\par
\noindent\hangindent=3.5em\hangafter=1\hspace*{3.5em}6.\ Mikusiński's Operational Calculus\nobreak\hfill\nobreak\leaders\hbox to 0.55em{\hss.\hss}\hfill\nobreak\ \pageref{sec:VI.6}\par
\noindent\hangindent=3.5em\hangafter=1\hspace*{3.5em}7.\ Sobolev's Lemma\nobreak\hfill\nobreak\leaders\hbox to 0.55em{\hss.\hss}\hfill\nobreak\ \pageref{sec:VI.7}\par
\noindent\hangindent=3.5em\hangafter=1\hspace*{3.5em}8.\ Gårding's Inequality\nobreak\hfill\nobreak\leaders\hbox to 0.55em{\hss.\hss}\hfill\nobreak\ \pageref{sec:VI.8}\par
\noindent\hangindent=3.5em\hangafter=1\hspace*{3.5em}9.\ Friedrichs' Theorem\nobreak\hfill\nobreak\leaders\hbox to 0.55em{\hss.\hss}\hfill\nobreak\ \pageref{sec:VI.9}\par
\noindent\hangindent=3.5em\hangafter=1\hspace*{3.5em}10.\ The Malgrange-Ehrenpreis Theorem\nobreak\hfill\nobreak\leaders\hbox to 0.55em{\hss.\hss}\hfill\nobreak\ \pageref{sec:VI.10}\par
\noindent\hangindent=3.5em\hangafter=1\hspace*{3.5em}11.\ Differential Operators with Uniform Strength\nobreak\hfill\nobreak\leaders\hbox to 0.55em{\hss.\hss}\hfill\nobreak\ \pageref{sec:VI.11}\par
\noindent\hangindent=3.5em\hangafter=1\hspace*{3.5em}12.\ The Hypoellipticity (Hörmänder's Theorem)\nobreak\hfill\nobreak\leaders\hbox to 0.55em{\hss.\hss}\hfill\nobreak\ \pageref{sec:VI.12}\par
\vspace{2.6mm}\noindent\hangindent=1.6em\hangafter=1\textbf{VII.\ \ Dual Operators}\nobreak\hfill\nobreak\leaders\hbox to 0.55em{\hss.\hss}\hfill\nobreak\ \pageref{chap:VII}\par
\noindent\hangindent=3.5em\hangafter=1\hspace*{3.5em}1.\ Dual Operators\nobreak\hfill\nobreak\leaders\hbox to 0.55em{\hss.\hss}\hfill\nobreak\ \pageref{sec:VII.1}\par
\noindent\hangindent=3.5em\hangafter=1\hspace*{3.5em}2.\ Adjoint Operators\nobreak\hfill\nobreak\leaders\hbox to 0.55em{\hss.\hss}\hfill\nobreak\ \pageref{sec:VII.2}\par
\noindent\hangindent=3.5em\hangafter=1\hspace*{3.5em}3.\ Symmetric Operators and Self-adjoint Operators\nobreak\hfill\nobreak\leaders\hbox to 0.55em{\hss.\hss}\hfill\nobreak\ \pageref{sec:VII.3}\par
\noindent\hangindent=3.5em\hangafter=1\hspace*{3.5em}4.\ Unitary Operators. The Cayley Transform\nobreak\hfill\nobreak\leaders\hbox to 0.55em{\hss.\hss}\hfill\nobreak\ \pageref{sec:VII.4}\par
\noindent\hangindent=3.5em\hangafter=1\hspace*{3.5em}5.\ The Closed Range Theorem\nobreak\hfill\nobreak\leaders\hbox to 0.55em{\hss.\hss}\hfill\nobreak\ \pageref{sec:VII.5}\par
\vspace{2.6mm}\noindent\hangindent=1.6em\hangafter=1\textbf{VIII.\ \ Resolvent and Spectrum}\nobreak\hfill\nobreak\leaders\hbox to 0.55em{\hss.\hss}\hfill\nobreak\ \pageref{chap:VIII}\par
\noindent\hangindent=3.5em\hangafter=1\hspace*{3.5em}1.\ The Resolvent and Spectrum\nobreak\hfill\nobreak\leaders\hbox to 0.55em{\hss.\hss}\hfill\nobreak\ \pageref{sec:VIII.1}\par
\noindent\hangindent=3.5em\hangafter=1\hspace*{3.5em}2.\ The Resolvent Equation and Spectral Radius\nobreak\hfill\nobreak\leaders\hbox to 0.55em{\hss.\hss}\hfill\nobreak\ \pageref{sec:VIII.2}\par
\noindent\hangindent=3.5em\hangafter=1\hspace*{3.5em}3.\ The Mean Ergodic Theorem\nobreak\hfill\nobreak\leaders\hbox to 0.55em{\hss.\hss}\hfill\nobreak\ \pageref{sec:VIII.3}\par
\noindent\hangindent=3.5em\hangafter=1\hspace*{3.5em}4.\ Ergodic Theorems of the Hille Type Concerning Pseudo-resolvents\nobreak\hfill\nobreak\leaders\hbox to 0.55em{\hss.\hss}\hfill\nobreak\ \pageref{sec:VIII.4}\par
\noindent\hangindent=3.5em\hangafter=1\hspace*{3.5em}5.\ The Mean Value of an Almost Periodic Function\nobreak\hfill\nobreak\leaders\hbox to 0.55em{\hss.\hss}\hfill\nobreak\ \pageref{sec:VIII.5}\par
\noindent\hangindent=3.5em\hangafter=1\hspace*{3.5em}6.\ The Resolvent of a Dual Operator\nobreak\hfill\nobreak\leaders\hbox to 0.55em{\hss.\hss}\hfill\nobreak\ \pageref{sec:VIII.6}\par
\noindent\hangindent=3.5em\hangafter=1\hspace*{3.5em}7.\ Dunford's Integral\nobreak\hfill\nobreak\leaders\hbox to 0.55em{\hss.\hss}\hfill\nobreak\ \pageref{sec:VIII.7}\par
\noindent\hangindent=3.5em\hangafter=1\hspace*{3.5em}8.\ The Isolated Singularities of a Resolvent\nobreak\hfill\nobreak\leaders\hbox to 0.55em{\hss.\hss}\hfill\nobreak\ \pageref{sec:VIII.8}\par
\vspace{2.6mm}\noindent\hangindent=1.6em\hangafter=1\textbf{IX.\ \ Analytical Theory of Semi-groups}\nobreak\hfill\nobreak\leaders\hbox to 0.55em{\hss.\hss}\hfill\nobreak\ \pageref{chap:IX}\par
\noindent\hangindent=3.5em\hangafter=1\hspace*{3.5em}1.\ The Semi-group of Class $(C_0)$\nobreak\hfill\nobreak\leaders\hbox to 0.55em{\hss.\hss}\hfill\nobreak\ \pageref{sec:IX.1}\par
\noindent\hangindent=3.5em\hangafter=1\hspace*{3.5em}2.\ The Equi-continuous Semi-group of Class $(C_0)$ in Locally Convex Spaces. Examples of Semi-groups\nobreak\hfill\nobreak\leaders\hbox to 0.55em{\hss.\hss}\hfill\nobreak\ \pageref{sec:IX.2}\par
\noindent\hangindent=3.5em\hangafter=1\hspace*{3.5em}3.\ The Infinitesimal Generator of an Equi-continuous Semi-group of Class $(C_0)$\nobreak\hfill\nobreak\leaders\hbox to 0.55em{\hss.\hss}\hfill\nobreak\ \pageref{sec:IX.3}\par
\noindent\hangindent=3.5em\hangafter=1\hspace*{3.5em}4.\ The Resolvent of the Infinitesimal Generator A\nobreak\hfill\nobreak\leaders\hbox to 0.55em{\hss.\hss}\hfill\nobreak\ \pageref{sec:IX.4}\par
\noindent\hangindent=3.5em\hangafter=1\hspace*{3.5em}5.\ Examples of Infinitesimal Generators\nobreak\hfill\nobreak\leaders\hbox to 0.55em{\hss.\hss}\hfill\nobreak\ \pageref{sec:IX.5}\par
\noindent\hangindent=3.5em\hangafter=1\hspace*{3.5em}6.\ The Exponential of a Continuous Linear Operator whose Powers are Equi-continuous\nobreak\hfill\nobreak\leaders\hbox to 0.55em{\hss.\hss}\hfill\nobreak\ \pageref{sec:IX.6}\par
\noindent\hangindent=3.5em\hangafter=1\hspace*{3.5em}7.\ The Representation and the Characterization of Equi-continuous Semi-groups of Class $(C_{0})$ in Terms of the Corresponding Infinitesimal Generators\nobreak\hfill\nobreak\leaders\hbox to 0.55em{\hss.\hss}\hfill\nobreak\ \pageref{sec:IX.7}\par
\noindent\hangindent=3.5em\hangafter=1\hspace*{3.5em}8.\ Contraction Semi-groups and Dissipative Operators\nobreak\hfill\nobreak\leaders\hbox to 0.55em{\hss.\hss}\hfill\nobreak\ \pageref{sec:IX.8}\par
\noindent\hangindent=3.5em\hangafter=1\hspace*{3.5em}9.\ Equi-continuous Groups of Class $(C_0)$. Stone's Theorem\nobreak\hfill\nobreak\leaders\hbox to 0.55em{\hss.\hss}\hfill\nobreak\ \pageref{sec:IX.9}\par
\noindent\hangindent=3.5em\hangafter=1\hspace*{3.5em}10.\ Holomorphic Semi-groups\nobreak\hfill\nobreak\leaders\hbox to 0.55em{\hss.\hss}\hfill\nobreak\ \pageref{sec:IX.10}\par
\noindent\hangindent=3.5em\hangafter=1\hspace*{3.5em}11.\ Fractional Powers of Closed Operators\nobreak\hfill\nobreak\leaders\hbox to 0.55em{\hss.\hss}\hfill\nobreak\ \pageref{sec:IX.11}\par
\noindent\hangindent=3.5em\hangafter=1\hspace*{3.5em}12.\ The Convergence of Semi-groups. The Trotter-Kato Theorem\nobreak\hfill\nobreak\leaders\hbox to 0.55em{\hss.\hss}\hfill\nobreak\ \pageref{sec:IX.12}\par
\noindent\hangindent=3.5em\hangafter=1\hspace*{3.5em}13.\ Dual Semi-groups. Phillips' Theorem\nobreak\hfill\nobreak\leaders\hbox to 0.55em{\hss.\hss}\hfill\nobreak\ \pageref{sec:IX.13}\par
\vspace{2.6mm}\noindent\hangindent=1.6em\hangafter=1\textbf{X.\ \ Compact Operators}\nobreak\hfill\nobreak\leaders\hbox to 0.55em{\hss.\hss}\hfill\nobreak\ \pageref{chap:X}\par
\noindent\hangindent=3.5em\hangafter=1\hspace*{3.5em}1.\ Compact Sets in $B$-spaces\nobreak\hfill\nobreak\leaders\hbox to 0.55em{\hss.\hss}\hfill\nobreak\ \pageref{sec:X.1}\par
\noindent\hangindent=3.5em\hangafter=1\hspace*{3.5em}2.\ Compact Operators and Nuclear Operators\nobreak\hfill\nobreak\leaders\hbox to 0.55em{\hss.\hss}\hfill\nobreak\ \pageref{sec:X.2}\par
\noindent\hangindent=3.5em\hangafter=1\hspace*{3.5em}3.\ The Rellich-Gårding Theorem\nobreak\hfill\nobreak\leaders\hbox to 0.55em{\hss.\hss}\hfill\nobreak\ \pageref{sec:X.3}\par
\noindent\hangindent=3.5em\hangafter=1\hspace*{3.5em}4.\ Schauder's Theorem\nobreak\hfill\nobreak\leaders\hbox to 0.55em{\hss.\hss}\hfill\nobreak\ \pageref{sec:X.4}\par
\noindent\hangindent=3.5em\hangafter=1\hspace*{3.5em}5.\ The Riesz-Schauder Theory\nobreak\hfill\nobreak\leaders\hbox to 0.55em{\hss.\hss}\hfill\nobreak\ \pageref{sec:X.5}\par
\noindent\hangindent=3.5em\hangafter=1\hspace*{3.5em}6.\ Dirichlet's Problem\nobreak\hfill\nobreak\leaders\hbox to 0.55em{\hss.\hss}\hfill\nobreak\ \pageref{sec:X.6}\par
\vspace{2.6mm}\noindent\hangindent=1.6em\hangafter=1\textbf{Appendix to Chapter X. The Nuclear Space of A. Grothendieck}\nobreak\hfill\nobreak\leaders\hbox to 0.55em{\hss.\hss}\hfill\nobreak\ \pageref{chap:AX}\par
\vspace{2.6mm}\noindent\hangindent=1.6em\hangafter=1\textbf{XI.\ \ Normed Rings and Spectral Representation}\nobreak\hfill\nobreak\leaders\hbox to 0.55em{\hss.\hss}\hfill\nobreak\ \pageref{chap:XI}\par
\noindent\hangindent=3.5em\hangafter=1\hspace*{3.5em}1.\ Maximal Ideals of a Normed Ring\nobreak\hfill\nobreak\leaders\hbox to 0.55em{\hss.\hss}\hfill\nobreak\ \pageref{sec:XI.1}\par
\noindent\hangindent=3.5em\hangafter=1\hspace*{3.5em}2.\ The Radical. The Semi-simplicity\nobreak\hfill\nobreak\leaders\hbox to 0.55em{\hss.\hss}\hfill\nobreak\ \pageref{sec:XI.2}\par
\noindent\hangindent=3.5em\hangafter=1\hspace*{3.5em}3.\ The Spectral Resolution of Bounded Normal Operators\nobreak\hfill\nobreak\leaders\hbox to 0.55em{\hss.\hss}\hfill\nobreak\ \pageref{sec:XI.3}\par
\noindent\hangindent=3.5em\hangafter=1\hspace*{3.5em}4.\ The Spectral Resolution of a Unitary Operator\nobreak\hfill\nobreak\leaders\hbox to 0.55em{\hss.\hss}\hfill\nobreak\ \pageref{sec:XI.4}\par
\noindent\hangindent=3.5em\hangafter=1\hspace*{3.5em}5.\ The Resolution of the Identity\nobreak\hfill\nobreak\leaders\hbox to 0.55em{\hss.\hss}\hfill\nobreak\ \pageref{sec:XI.5}\par
\noindent\hangindent=3.5em\hangafter=1\hspace*{3.5em}6.\ The Spectral Resolution of a Self-adjoint Operator\nobreak\hfill\nobreak\leaders\hbox to 0.55em{\hss.\hss}\hfill\nobreak\ \pageref{sec:XI.6}\par
\noindent\hangindent=3.5em\hangafter=1\hspace*{3.5em}7.\ Real Operators and Semi-bounded Operators. Friedrichs' Theorem\nobreak\hfill\nobreak\leaders\hbox to 0.55em{\hss.\hss}\hfill\nobreak\ \pageref{sec:XI.7}\par
\noindent\hangindent=3.5em\hangafter=1\hspace*{3.5em}8.\ The Spectrum of a Self-adjoint Operator. Rayleigh's Principle and the Krylov-Weinstein Theorem. The Multiplicity of the Spectrum\nobreak\hfill\nobreak\leaders\hbox to 0.55em{\hss.\hss}\hfill\nobreak\ \pageref{sec:XI.8}\par
\noindent\hangindent=3.5em\hangafter=1\hspace*{3.5em}9.\ The General Expansion Theorem. A Condition for the Absence of the Continuous Spectrum\nobreak\hfill\nobreak\leaders\hbox to 0.55em{\hss.\hss}\hfill\nobreak\ \pageref{sec:XI.9}\par
\noindent\hangindent=3.5em\hangafter=1\hspace*{3.5em}10.\ The Peter-Weyl-Neumann Theory\nobreak\hfill\nobreak\leaders\hbox to 0.55em{\hss.\hss}\hfill\nobreak\ \pageref{sec:XI.10}\par
\noindent\hangindent=3.5em\hangafter=1\hspace*{3.5em}11.\ Tannaka's Duality Theorem for Non-commutative Compact Groups\nobreak\hfill\nobreak\leaders\hbox to 0.55em{\hss.\hss}\hfill\nobreak\ \pageref{sec:XI.11}\par
\noindent\hangindent=3.5em\hangafter=1\hspace*{3.5em}12.\ Functions of a Self-adjoint Operator\nobreak\hfill\nobreak\leaders\hbox to 0.55em{\hss.\hss}\hfill\nobreak\ \pageref{sec:XI.12}\par
\noindent\hangindent=3.5em\hangafter=1\hspace*{3.5em}13.\ Stone's Theorem and Bochner's Theorem\nobreak\hfill\nobreak\leaders\hbox to 0.55em{\hss.\hss}\hfill\nobreak\ \pageref{sec:XI.13}\par
\noindent\hangindent=3.5em\hangafter=1\hspace*{3.5em}14.\ A Canonical Form of a Self-adjoint Operator with Simple Spectrum\nobreak\hfill\nobreak\leaders\hbox to 0.55em{\hss.\hss}\hfill\nobreak\ \pageref{sec:XI.14}\par
\noindent\hangindent=3.5em\hangafter=1\hspace*{3.5em}15.\ The Defect Indices of a Symmetric Operator. The Generalized Resolution of the Identity\nobreak\hfill\nobreak\leaders\hbox to 0.55em{\hss.\hss}\hfill\nobreak\ \pageref{sec:XI.15}\par
\noindent\hangindent=3.5em\hangafter=1\hspace*{3.5em}16.\ The Group-ring $L^1$ and Wiener's Tauberian Theorem\nobreak\hfill\nobreak\leaders\hbox to 0.55em{\hss.\hss}\hfill\nobreak\ \pageref{sec:XI.16}\par
\vspace{2.6mm}\noindent\hangindent=1.6em\hangafter=1\textbf{XII.\ \ Other Representation Theorems in Linear Spaces}\nobreak\hfill\nobreak\leaders\hbox to 0.55em{\hss.\hss}\hfill\nobreak\ \pageref{chap:XII}\par
\noindent\hangindent=3.5em\hangafter=1\hspace*{3.5em}1.\ Extremal Points. The Krein-Milman Theorem\nobreak\hfill\nobreak\leaders\hbox to 0.55em{\hss.\hss}\hfill\nobreak\ \pageref{sec:XII.1}\par
\noindent\hangindent=3.5em\hangafter=1\hspace*{3.5em}2.\ Vector Lattices\nobreak\hfill\nobreak\leaders\hbox to 0.55em{\hss.\hss}\hfill\nobreak\ \pageref{sec:XII.2}\par
\noindent\hangindent=3.5em\hangafter=1\hspace*{3.5em}3.\ $B$-lattices and $F$-lattices\nobreak\hfill\nobreak\leaders\hbox to 0.55em{\hss.\hss}\hfill\nobreak\ \pageref{sec:XII.3}\par
\noindent\hangindent=3.5em\hangafter=1\hspace*{3.5em}4.\ A Convergence Theorem of Banach\nobreak\hfill\nobreak\leaders\hbox to 0.55em{\hss.\hss}\hfill\nobreak\ \pageref{sec:XII.4}\par
\noindent\hangindent=3.5em\hangafter=1\hspace*{3.5em}5.\ The Representation of a Vector Lattice as Point Functions\nobreak\hfill\nobreak\leaders\hbox to 0.55em{\hss.\hss}\hfill\nobreak\ \pageref{sec:XII.5}\par
\noindent\hangindent=3.5em\hangafter=1\hspace*{3.5em}6.\ The Representation of a Vector Lattice as Set Functions\nobreak\hfill\nobreak\leaders\hbox to 0.55em{\hss.\hss}\hfill\nobreak\ \pageref{sec:XII.6}\par
\vspace{2.6mm}\noindent\hangindent=1.6em\hangafter=1\textbf{XIII.\ \ Ergodic Theory and Diffusion Theory}\nobreak\hfill\nobreak\leaders\hbox to 0.55em{\hss.\hss}\hfill\nobreak\ \pageref{chap:XIII}\par
\noindent\hangindent=3.5em\hangafter=1\hspace*{3.5em}1.\ The Markov Process with an Invariant Measure*\nobreak\hfill\nobreak\leaders\hbox to 0.55em{\hss.\hss}\hfill\nobreak\ \pageref{sec:XIII.1}\par
\noindent\hangindent=3.5em\hangafter=1\hspace*{3.5em}2.\ An Individual Ergodic Theorem and Its Applications\nobreak\hfill\nobreak\leaders\hbox to 0.55em{\hss.\hss}\hfill\nobreak\ \pageref{sec:XIII.2}\par
\noindent\hangindent=3.5em\hangafter=1\hspace*{3.5em}3.\ The Ergodic Hypothesis and the $H$-theorem\nobreak\hfill\nobreak\leaders\hbox to 0.55em{\hss.\hss}\hfill\nobreak\ \pageref{sec:XIII.3}\par
\noindent\hangindent=3.5em\hangafter=1\hspace*{3.5em}4.\ The Ergodic Decomposition of a Markov Process with a Locally Compact Phase Space\nobreak\hfill\nobreak\leaders\hbox to 0.55em{\hss.\hss}\hfill\nobreak\ \pageref{sec:XIII.4}\par
\noindent\hangindent=3.5em\hangafter=1\hspace*{3.5em}5.\ The Brownian Motion on a Homogeneous Riemannian Space\nobreak\hfill\nobreak\leaders\hbox to 0.55em{\hss.\hss}\hfill\nobreak\ \pageref{sec:XIII.5}\par
\noindent\hangindent=3.5em\hangafter=1\hspace*{3.5em}6.\ The Generalized Laplacian of W. Feller\nobreak\hfill\nobreak\leaders\hbox to 0.55em{\hss.\hss}\hfill\nobreak\ \pageref{sec:XIII.6}\par
\noindent\hangindent=3.5em\hangafter=1\hspace*{3.5em}7.\ An Extension of the Diffusion Operator\nobreak\hfill\nobreak\leaders\hbox to 0.55em{\hss.\hss}\hfill\nobreak\ \pageref{sec:XIII.7}\par
\noindent\hangindent=3.5em\hangafter=1\hspace*{3.5em}8.\ Markov Processes and Potentials\nobreak\hfill\nobreak\leaders\hbox to 0.55em{\hss.\hss}\hfill\nobreak\ \pageref{sec:XIII.8}\par
\noindent\hangindent=3.5em\hangafter=1\hspace*{3.5em}9.\ Abstract Potential Operators and Semi-groups\nobreak\hfill\nobreak\leaders\hbox to 0.55em{\hss.\hss}\hfill\nobreak\ \pageref{sec:XIII.9}\par
\vspace{2.6mm}\noindent\hangindent=1.6em\hangafter=1\textbf{XIV.\ \ The Integration of the Equation of Evolution*}\nobreak\hfill\nobreak\leaders\hbox to 0.55em{\hss.\hss}\hfill\nobreak\ \pageref{chap:XIV}\par
\noindent\hangindent=3.5em\hangafter=1\hspace*{3.5em}1.\ Integration of Diffusion Equations in $L^{2}(\mathbb{R}^{m})$\nobreak\hfill\nobreak\leaders\hbox to 0.55em{\hss.\hss}\hfill\nobreak\ \pageref{sec:XIV.1}\par
\noindent\hangindent=3.5em\hangafter=1\hspace*{3.5em}2.\ Integration of Diffusion Equations in a Compact Riemannian Space\nobreak\hfill\nobreak\leaders\hbox to 0.55em{\hss.\hss}\hfill\nobreak\ \pageref{sec:XIV.2}\par
\noindent\hangindent=3.5em\hangafter=1\hspace*{3.5em}3.\ Integration of Wave Equations in a Euclidean Space $\mathbb{R}^{m}$\nobreak\hfill\nobreak\leaders\hbox to 0.55em{\hss.\hss}\hfill\nobreak\ \pageref{sec:XIV.3}\par
\noindent\hangindent=3.5em\hangafter=1\hspace*{3.5em}4.\ Integration of Temporally Inhomogeneous Equations of Evolution in a $B$-space\nobreak\hfill\nobreak\leaders\hbox to 0.55em{\hss.\hss}\hfill\nobreak\ \pageref{sec:XIV.4}\par
\noindent\hangindent=3.5em\hangafter=1\hspace*{3.5em}5.\ The Method of Tanabe and Sobolevski\nobreak\hfill\nobreak\leaders\hbox to 0.55em{\hss.\hss}\hfill\nobreak\ \pageref{sec:XIV.5}\par
\noindent\hangindent=3.5em\hangafter=1\hspace*{3.5em}6.\ Non-linear Evolution Equations 1 (The Kōmura-Kato Approach)\nobreak\hfill\nobreak\leaders\hbox to 0.55em{\hss.\hss}\hfill\nobreak\ \pageref{sec:XIV.6}\par
\noindent\hangindent=3.5em\hangafter=1\hspace*{3.5em}7.\ Non-linear Evolution Equations 2 (The Approach through the Crandall-Liggett Convergence Theorem)\nobreak\hfill\nobreak\leaders\hbox to 0.55em{\hss.\hss}\hfill\nobreak\ \pageref{sec:XIV.7}\par
\vspace{2.6mm}\noindent\hangindent=1.6em\hangafter=1\textbf{Bibliography}\nobreak\hfill\nobreak\leaders\hbox to 0.55em{\hss.\hss}\hfill\nobreak\ \pageref{chap:Bibliography}\par
\vspace{2.6mm}\noindent\hangindent=1.6em\hangafter=1\textbf{Index}\nobreak\hfill\nobreak\leaders\hbox to 0.55em{\hss.\hss}\hfill\nobreak\ \pageref{chap:Index}\par
